% LANA_SELECTION_DEFINITION_CLEANUP_20260708T182812Z
% Paper: M1 RP-2 / m1_rp2_environment_quenching
% Lane-local revision draft only; does not overwrite the current linked manuscript or PDF.
\documentclass[twocolumn]{aastex631}
\usepackage{amsmath}
\usepackage{booktabs}
\graphicspath{{/Users/duhokim/NebulaMind/NebulaMind/.hermes/handoffs/galaxy-evolution/mastermind/aas-autopilot/runs/SDSS_REMAINING_TOPIC_PILOTS_20260708T125828Z/m1_rp2_environment_quenching/figures/}}

\shorttitle{Selection-Flagged Density-Proxy Quenching}
\shortauthors{NebulaMind Autopilot}

\begin{document}

\title{A Selection-Flagged Nearest-Neighbour Density Proxy for Quenching in an SDSS DR17 Emission-Line Pilot Sample}

\author{NebulaMind Research Autopilot}
\affiliation{Local reproducible pilot run; public SDSS DR17 data only}

\begin{abstract}
We revise the M1 RP-2 pilot into a bounded, selection-flagged SDSS density-proxy paper.  The active measurement is not a halo-environment or central/satellite test; it asks whether an internally computed 10th-nearest-neighbour density ranking is associated with the quenched fraction inside a capped SDSS DR17 four-line emission sample.  Public SDSS count checks show 249,917 galaxies satisfying the strict four-BPT-line S/N$\geq3$ selection, while the cached pilot contains the first 60,000 rows returned by the preserved \texttt{TOP 60000 ... ORDER BY specObjID} query, i.e. 24.0\% of that strict parent.  Quenching is now explicitly defined as catalog $\log({\rm sSFR}/{\rm yr}^{-1})< -11.0$.  Within the cached denominator, the high-density quartile contains 3,456 quenched galaxies out of 15,000 ($f_Q=0.230$), while the low-density quartile contains 2,710 out of 15,000 ($f_Q=0.181$).  The high-minus-low difference is 0.049 with bootstrap 95\% interval [0.041, 0.059], and a linear-probability diagnostic adjusted for stellar mass and redshift gives $0.032\pm0.004$.  The result is a reproducible local-density association, not causal proof of environmental quenching.
\end{abstract}

\keywords{galaxies: evolution --- galaxies: environments --- galaxies: star formation --- surveys --- methods: data analysis}

\section{Purpose and Scope}\label{sec:scope}
Mass and environment are long-standing organizing axes for galaxy quenching \citep{peng2010,baldry2006}.  This pilot deliberately tests only the part that the preserved SDSS data can support: a within-sample density ranking versus a catalog-sSFR quenching flag.  It does not include group catalogues, halo masses, satellite infall histories, morphology, or forward-modeled survey boundaries, all of which are needed before interpreting the association as environmental quenching in the stronger sense \citep{wetzel2013,goubert2024}.

\section{Data, Selection Function, and Operational Definitions}\label{sec:data}
The underlying measurements come from public SDSS DR17 spectroscopy and value-added catalog quantities \citep{abdurrouf2022,york2000,brinchmann2004}.  The local run uses galaxies with spectroscopic class \texttt{GALAXY}, $0.02<z<0.12$, finite catalog stellar mass and sSFR, and positive H$\alpha$, H$\beta$, [O~III]$\lambda5007$, and [N~II]$\lambda6584$ fluxes with S/N$\geq3$ in all four lines.  BPT line-ratio classes are recomputed with the standard Kauffmann and Kewley demarcations \citep{baldwin1981,kauffmann2003,kewley2001}, but the density result below depends on the quenching flag rather than the optical-AGN class.

Table~\ref{tab:rp2_selection} is mandatory context for any incidence wording.  The cached 60,000-row table is a capped, ordered subset of the strict four-line SDSS parent, not a random or complete local-galaxy sample.  Four-line retention is also sSFR-dependent: the public SDSS count check retains 32,021 of 95,424 galaxies (33.56\%) in the $-12<\log{\rm sSFR}<-11$ bin, but 84,488 of 89,075 (94.85\%) in the $-10<\log{\rm sSFR}<-9.5$ bin.  That differential retention is especially relevant for a quenching paper.

\begin{deluxetable*}{lrrl}
\tablecaption{Selection-function disclosure for the shared SDSS four-line denominator\label{tab:rp2_selection}}
\tablehead{\colhead{Selection stage} & \colhead{Public SDSS DR17 count} & \colhead{Cached count used here} & \colhead{Manuscript implication}}
\startdata
Spectroscopic \texttt{GALAXY}, $0.02<z<0.12$ & 501,060 & --- & Redshift-window parent only \\
Joined catalog with mass/sSFR bounds & 416,554 & --- & Value-added-property parent \\
Positive four BPT fluxes and errors & 373,445 & 60,000 & Four-line emission denominator begins to exclude weak-line systems \\
Four BPT lines with S/N$\geq3$ & 249,917 & 60,000 & Cached sample covers 24.0\% of strict parent \\
Four BPT lines with S/N$\geq5$ & 176,523 & 42,446 & Robustness subset, not the headline denominator \\
Four BPT lines with S/N$\geq10$ & 91,768 & 22,311 & High-S/N subset changes population mix \\
\enddata
\tablecomments{Counts are inherited from the read-only Tori/Goru selection-function packets.  The cached query used \texttt{TOP 60000 ... ORDER BY specObjID}; therefore all fractions in this draft are conditional on the cached denominator.}
\end{deluxetable*}

The quenching flag is now explicit: a galaxy is counted as quenched when \texttt{specsfr\_tot\_p50}$<-11.0$.  The density observable is likewise explicit: approximate low-redshift comoving Cartesian positions are built from $(\alpha,\delta,z)$, the distance to the 10th nearest neighbour is measured, and the local-density proxy is proportional to $10/(4\pi d_{10}^{3}/3)$.  Quartiles are computed inside the cached sample.  Because survey masks, edge corrections, fiber-collision corrections, volume completeness, and group membership are not modeled, this is a rank-order density proxy only.

\section{Results}\label{sec:results}
The high-density quartile has a larger quenched fraction than the low-density quartile in the cached pilot measurement.  Table~\ref{tab:rp2_summary} replaces the earlier itemized result with explicit denominators, threshold language, and the interval convention.

\begin{deluxetable*}{lccc}
\tablecaption{Density-proxy quenching summary for the cached SDSS DR17 pilot\label{tab:rp2_summary}}
\tablehead{\colhead{Quantity} & \colhead{Low-density quartile} & \colhead{High-density quartile} & \colhead{Contrast or model term}}
\startdata
Galaxies in quartile & 15,000 & 15,000 & --- \\
Quenched definition & \multicolumn{3}{c}{\texttt{specsfr\_tot\_p50}$<-11.0$} \\
Quenched count & 2,710 & 3,456 & --- \\
Quenched fraction & 0.181 $\pm$ 0.003 & 0.230 $\pm$ 0.003 & $\Delta f_Q=0.049$ \\
Bootstrap interval for $\Delta f_Q$ & --- & --- & [0.041, 0.059] \\
Mass--redshift adjusted LPM coefficient & --- & --- & $0.032\pm0.004$ \\
Approximate 95\% interval for LPM coefficient & --- & --- & [0.025, 0.040] \\
\enddata
\tablecomments{Fractions and bootstrap intervals are copied from the preserved topic-specific JSON.  The linear-probability-model coefficient is a diagnostic adjusted for stellar mass and redshift; it is not a causal environmental-quenching estimate.}
\end{deluxetable*}

\begin{figure}
\centering
\includegraphics[width=\columnwidth]{m1\_rp2\_environment\_quenching\_figure1.pdf}
\caption{Preserved topic-specific SDSS DR17 density-proxy diagnostic.  The plotted comparison should be interpreted as a cached-sample nearest-neighbour rank test with the quenching threshold in Table~\ref{tab:rp2_summary}, not as a calibrated halo-environment measurement.}
\label{fig:rp2_density}
\end{figure}

\section{Discussion and Integration Notes}\label{sec:discussion}
The result supports a cautious statement: within this cached emission-line denominator, the high-density quartile is more often low-sSFR than the low-density quartile.  It does not isolate the physical mechanism.  Morphology--density correlations, stellar age, halo mass, satellite status, aperture scale, and the sSFR-dependent line-selection function can all contribute.  The appropriate next integration step is to keep Peng/Baldry-style mass--environment context near the motivation, then reserve group-catalogue and simulation-comparison citations for the missing-data paragraph rather than for the measured result itself.

\section{Conclusions}\label{sec:conclusions}
\begin{enumerate}
\item The revised M1 RP-2 draft now names the key operational thresholds: four BPT lines with S/N$\geq3$, a cached \texttt{TOP 60000} denominator, \texttt{specsfr\_tot\_p50}$<-11.0$ for quenching, and a 10th-nearest-neighbour density rank.
\item The measured high-minus-low density-quartile quenched-fraction difference is 0.049 with bootstrap 95\% interval [0.041, 0.059].
\item Selection disclosure is central: the cached sample covers 24.0\% of the strict public S/N$\geq3$ parent and four-line retention is strongly sSFR-dependent.
\item The paper should remain an SDSS density-proxy association baseline, not a causal proof of environmental quenching.
\end{enumerate}

\section*{Reproducibility and Safety Note}
Source analysis summary: \texttt{m1\_rp2\_environment\_quenching/analysis\_results.json} under run SDSS\_REMAINING\_TOPIC\_PILOTS\_20260708T125828Z.  This Lana revision is lane-local and did not overwrite the public-linked manuscript or PDF.

\acknowledgments
This pilot used public SDSS DR17 data products and open-source Python tools.

\begin{thebibliography}{}
\bibitem[Abdurro'uf et al.(2022)]{abdurrouf2022} Abdurro'uf, A., et al. 2022, ApJS, 259, 35
\bibitem[Baldry et al.(2006)]{baldry2006} Baldry, I.~K., et al. 2006, MNRAS, 373, 469
\bibitem[Baldwin et al.(1981)]{baldwin1981} Baldwin, J.~A., Phillips, M.~M., \& Terlevich, R. 1981, PASP, 93, 5
\bibitem[Brinchmann et al.(2004)]{brinchmann2004} Brinchmann, J., Charlot, S., White, S.~D.~M., et al. 2004, MNRAS, 351, 1151
\bibitem[Goubert et al.(2024)]{goubert2024} Goubert, P.~H., et al. 2024, arXiv:2401.12953
\bibitem[Kauffmann et al.(2003)]{kauffmann2003} Kauffmann, G., Heckman, T.~M., Tremonti, C., et al. 2003, MNRAS, 346, 1055
\bibitem[Kewley et al.(2001)]{kewley2001} Kewley, L.~J., Dopita, M.~A., Sutherland, R.~S., Heisler, C.~A., \& Trevena, J. 2001, ApJ, 556, 121
\bibitem[Peng et al.(2010)]{peng2010} Peng, Y.-j., Lilly, S.~J., Kova\v{c}, K., et al. 2010, ApJ, 721, 193
\bibitem[Wetzel et al.(2013)]{wetzel2013} Wetzel, A.~R., Tinker, J.~L., Conroy, C., \& van den Bosch, F.~C. 2013, MNRAS, 432, 336
\bibitem[York et al.(2000)]{york2000} York, D.~G., Adelman, J., Anderson, J.~E., Jr., et al. 2000, AJ, 120, 1579
\end{thebibliography}

\end{document}
